\documentclass[working, oneside]{inputs/tuftebook}
%%%%%%%%%%%%%%%%%%%%
%% SUPER PREAMBLE %%
%%%%%%%%%%%%%%%%%%%%

\documentclass[a4paper]{article}
\usepackage[utf8]{inputenc}
\usepackage[T1]{fontenc} % Fonts and stuff
\usepackage{amsmath, amsfonts, mathtools, amsthm, amssymb} % math

\usepackage{fancyhdr} % Header, Footer etc.
\usepackage{adforn}

\pagestyle{fancy}
\fancyhf{}
\fancyhead[R]{Albert Lunde, John Faxe Jensen}
\fancyfoot[C]{\adforn{17}\quad\thepage\quad\adforn{45}}

\renewcommand{\headrule}{%
	\hrulefill
	 {\quad\adforn{21}\adforn{10}\adforn{49}\quad}%
	\hrulefill}
\renewcommand{\footrulewidth}{0pt}

%% Margin Control %%

\def\changemargin#1#2{\list{}{\rightmargin#2\leftmargin#1}\item[]}
\let\endchangemargin=\endlist


%%%%%%%%%%%%%%%%%%%%%%%%%%%%%%%%%%%%%%%%%%%%%%%%%%%%%%%%%%

% figure support

\usepackage{import}
\usepackage{pdfpages}
\usepackage{transparent}
\usepackage{xcolor}

\newcommand{\incfig}[2][1]{%
    \def\svgwidth{#1\columnwidth}
    \import{figures/}{#2.pdf_tex}
}

\pdfsuppresswarningpagegroup=1

%%%%%%%%%%%%%%%%%%%%%%%%%%%%%%%%%%%%%%%%%%%%%%%%%%%%%%%%%%

\usepackage{tikzsymbols} % Symbols
\usepackage{mdframed} % Boxes around theorem environments
\usepackage{thmtools}


% Exercise environment with surrounding box

\declaretheoremstyle[
    name=Exercise,
    mdframed={
  skipabove=0pt,
  skipbelow=20pt,
  innerleftmargin=10pt,
  innerrightmargin=10pt,
  innerbottommargin=7pt}
]{thmsty}
\declaretheorem[style=thmsty,numbered=no]{exercise}

% Solution environment, with coffee cup

\newenvironment{solution}
 {\renewcommand\qedsymbol{\tikzsymbolsuse{Coffeecup}}\begin{proof}[Solution]}
 {\end{proof}}

% Subexercise enviroment
%  \newenvironment{subexercise}[1]
%  {
% 	\begin{changemargin}{1.0cm}{1.0cm}
% 	\textbf{(#1)}\em
% 	}{
	% \end{changemargin}
	% } 

 \newenvironment{subexercise}[1]
 {\noindent
	 \textbf{(#1)} \quad \adforn{10} \quad \em
 }{}

% Mathematical typesetting stuff.

 \newcommand{\dd}{\mathrm{\textbf{d}}}



\usepackage{pdfpages}

\usepackage{lipsum}
\usepackage{parskip}
\usepackage{titletoc}

\usepackage{cmbright}
\usepackage{bm}
\addbibresource{bibliography.bib}

\begin{document}
% \includepdf[pages=-, fitpaper=true]{inputs/frontmatter.pdf}
\section*{Introduction}
\section*{Theory}
\subsection*{Spectrometer}
We shall first, in short, explain how a spectrometer works and what it is used. A spectrometer is a device that can measure the spectral composition of light. Fig. 1 shows the spectrum of hydrogen. The most common form of spectrometer, which is also what we have used, is a grating spectrometer. A grating spectrometer, separates the light by wavelength with a periodic grate, as illustrated in fig 2. Let us a consider a monochromatic component of light. As it is reflected off the grate it will only create construcive interference at specific angles. This angle is related to the wavelength by the follwing formula called the 'grating equation'
\[
d\left( \sin \theta_i - \sin \theta_m \right) = \lambda m \quad \text{where}\quad  m = \pm 0, \pm 1, \pm 2 \ldots
.\] 
\subsection*{Absorbtion}
\subsection*{Emission}
All elements have different spectra. These are to the specific energy levels where they are able to absorb light (different wavelengths of light have different energy levels). These energy levels are the possible excited states that the elements electronic structure can be in. An element emits light, when its electron structure descends in energy. The energy of light emitted is then equal to the magnitude of the descent. As a consequence an element can only emit light with specific energies.
\section*{Results}
\section*{Discussion}
\section*{Conclusion}
\end{document}
